% Layout fixture: a deck with every kind of overlay step and a header and footer bar theme.
% \pause and \item<n-> add to the earlier steps (an earlier page is all on the last one);
% \only<1>{..}\only<2>{..} replaces text (the earlier page has text the last one lacks);
% \uncover shows more of the same slide. Written for ReadableSTEM's tests (MIT).
\documentclass{beamer}
\usepackage{amsmath}
\usetheme{Warsaw}
\setbeamertemplate{navigation symbols}{}

\title{Waves and Oscillations}
\author{R. Example}

\begin{document}

\section{Waves}
\begin{frame}{Wave Speed}
  A wave of wavelength $\lambda$ and frequency $f$ travels at
  \pause
  \begin{equation*}
    v = \lambda f.
  \end{equation*}
  \pause
  Longer waves move faster at a fixed frequency.
\end{frame}

\begin{frame}{Cases}
  \only<1>{Light waves need no medium.}%
  \only<2>{Sound waves need a medium.}
\end{frame}

\begin{frame}{Checks}
  \begin{itemize}
    \item<1-> The period is $T = 1/f$.
    \item<2-> The angular frequency is $\omega = 2\pi f$.
    \item<3-> The wave number is $k = 2\pi/\lambda$.
    \item<4-> Speed is $v = \omega/k$.
  \end{itemize}
\end{frame}

\section{Superposition}
\begin{frame}{Uncovered}
  \uncover<1->{Two waves add.}

  \uncover<2->{Where crests meet, the sum is larger.}

  \uncover<3->{Where a crest meets a trough, it cancels.}
\end{frame}

\begin{frame}{Summary}
  Frequency, wavelength and speed are linked by $v = \lambda f$.
\end{frame}

\end{document}
